\chapter{Methodology}
\label{chapter2}

<Everything that comes under the `Methodology' criterion in the mark scheme should be described in one, or possibly more than one, chapter(s).>

\section{Requirements and design goals}

The literature reviewed in chapter 1 identifies a gap in the comparative empirical evaluation of management polices for ASP-mediated UAS Traffic management (UTM), closing that gap requires a simualtion framework that is sufficialtly faithful to hte operational context of UTM to make its results meaningful, and sufficiently flexible to support controlled comparison across a range of polcies and operating conditons. This section sets out the requirements created with these considerations in mine. which together infrom the design choices documented in teh remainder of the chapter. 

The functional requirements are as follows. The framework must model an urban airspace as a network of independent Airspace Service providers (ASPs), each controlling thier defined area and handling the handover of traffic to its neighbouring controllers. It must process flight plan requests dynamically during execution rather then as a fixed batch known in advance (to reflect the real world) it must support multiple management from MAPF reviewed in the chapter above. it must permit policies to be tested in the same envrionments and configurations. it must record sufficent data on each flight and simuaitno to support teh saftey, efficency and coordination metrics described. it must support an optional tacticle collision avoidance layer, drawn from the velocity obstacle family reviewed in the reservation section, that can be applied above any policy. 

the non functional requries are equally importatnt. The frameowrk must be deteminsitic given a fixed random seed, so that experimental results are reporducible. it must be modular enough to permit new polices and new tactical layers to be added without modifying the siulation core, in support of both the comparitive evaluations undertaken in this work and any future extensions. It must provid  configeration through external files rather then paremeters so experimental conditions can be varied without recompliation. it must complete experimetns within a resanable time on comodity hardware as it will be ran many time in many configeration.

From these requirements they derive several other design goals that recur throughout the rest of this chapter. The simulation must be discrete event driven rather than time-stepped. Many considerations motivated this choice. First, performance: a time-stepped simulation processes every agent (UAV) at every clock tick regardless of whether there state has changed, while a discrete-event simulation processes a UAV only when an event affects it, which scales considerably better as the number of UAVs and the length of the simulated horizon grow. Secondly, precision in collision detection: continuous time geometric collision detection over analytic trajectory segments yields the exact moment and distance of closest approach, while a fixed-tick simulation can miss close-approach events that occur entirely within a single time step unless the step is made impractically small. Third, reproducibility: a discrete-event simulation is fully determined by its random seed and configuration, whereas a time-stepped simulation introduces the time step itself as a hidden parameter whose value can affect experimental results. SimPy \cite{simpy2024} is the natural choice for this paradigm in Python.

These requirements along with design goals have driven the architecture presented in <section>, the simulation model will be described in <section> and the policy abstraction in <section>.

There will also be a section on the revervation system <section> and a summary on how the project was managed <section>
